\documentclass[10pt,a4paper]{amsart}
\usepackage{amsmath,amssymb,amsthm,graphicx}
\usepackage{verbatim}
\newenvironment{reference}{\comment}{\endcomment}
%\newenvironment{reference}{}{}
\newenvironment{slogan}{\comment}{\endcomment}
\newenvironment{history}{\comment}{\endcomment}

% For commutative diagrams we use Xy-pic
\usepackage[all]{xy}

% We use 2cell for 2-commutative diagrams.
\xyoption{2cell}
% Theorem environments.
%
\theoremstyle{plain}
\newtheorem{theorem}[subsection]{Theorem}
\newtheorem{proposition}[subsection]{Proposition}
\newtheorem{lemma}[subsection]{Lemma}

\theoremstyle{definition}
\newtheorem{definition}[subsection]{Definition}
\newtheorem{example}[subsection]{Example}
\newtheorem{exercise}[subsection]{Exercise}
\newtheorem{situation}[subsection]{Situation}

\theoremstyle{remark}
\newtheorem{remark}[subsection]{Remark}
\newtheorem{remarks}[subsection]{Remarks}

\numberwithin{equation}{subsection}

% Macros
%
\def\lim{\mathop{\mathrm{lim}}\nolimits}
\def\colim{\mathop{\mathrm{colim}}\nolimits}
\def\Spec{\mathop{\mathrm{Spec}}}
\def\Hom{\mathop{\mathrm{Hom}}\nolimits}
\def\Ext{\mathop{\mathrm{Ext}}\nolimits}
\def\SheafHom{\mathop{\mathcal{H}\!\mathit{om}}\nolimits}
\def\SheafExt{\mathop{\mathcal{E}\!\mathit{xt}}\nolimits}
\def\Sch{\mathit{Sch}}
\def\Mor{\mathop{\mathrm{Mor}}\nolimits}
\def\Ob{\mathop{\mathrm{Ob}}\nolimits}
\def\Sh{\mathop{\mathit{Sh}}\nolimits}
\def\NL{\mathop{N\!L}\nolimits}
\def\CH{\mathop{\mathrm{CH}}\nolimits}
\def\proetale{{pro\text{-}\acute{e}tale}}
\def\etale{{\acute{e}tale}}
\def\QCoh{\mathit{QCoh}}
\def\Ker{\mathop{\mathrm{Ker}}}
\def\Im{\mathop{\mathrm{Im}}}
\def\Coker{\mathop{\mathrm{Coker}}}
\def\Coim{\mathop{\mathrm{Coim}}}

% Boxtimes
%
\DeclareMathSymbol{\boxtimes}{\mathbin}{AMSa}{"02}

%
% Macros for moduli stacks/spaces
%
\def\QCohstack{\mathcal{QC}\!\mathit{oh}}
\def\Cohstack{\mathcal{C}\!\mathit{oh}}
\def\Spacesstack{\mathcal{S}\!\mathit{paces}}
\def\Quotfunctor{\mathrm{Quot}}
\def\Hilbfunctor{\mathrm{Hilb}}
\def\Curvesstack{\mathcal{C}\!\mathit{urves}}
\def\Polarizedstack{\mathcal{P}\!\mathit{olarized}}
\def\Complexesstack{\mathcal{C}\!\mathit{omplexes}}
% \Pic is the operator that assigns to X its picard group, usage \Pic(X)
% \Picardstack_{X/B} denotes the Picard stack of X over B
% \Picardfunctor_{X/B} denotes the Picard functor of X over B
\def\Pic{\mathop{\mathrm{Pic}}\nolimits}
\def\Picardstack{\mathcal{P}\!\mathit{ic}}
\def\Picardfunctor{\mathrm{Pic}}
\def\Deformationcategory{\mathcal{D}\!\mathit{ef}}
\usepackage{cleveref}
\crefname{originallemma}{original lemma}{original lemmas}
\Crefname{originallemma}{Original lemma}{Original lemmas}
\makeatletter
\newlabel{lemma-quasi-compact-local}{{03KG}{source}{}{sourcepointer-1}{}}
\newlabel{lemma-quasi-compact-local@cref}{{[originallemma][0][]03KG}{[1][0][]source}{}{}{}}
\makeatother
\crefname{originallemma}{original lemma}{original lemmas}
\Crefname{originallemma}{Original lemma}{Original lemmas}
\makeatletter
\newlabel{lemma-quasi-compact-permanence}{{03KS}{source}{}{sourcepointer-2}{}}
\newlabel{lemma-quasi-compact-permanence@cref}{{[originallemma][0][]03KS}{[1][0][]source}{}{}{}}
\makeatother
\crefname{originallemma}{original lemma}{original lemmas}
\Crefname{originallemma}{Original lemma}{Original lemmas}
\makeatletter
\newlabel{morphisms-lemma-universally-closed-quasi-compact}{{04XU}{source}{}{sourcepointer-3}{}}
\newlabel{morphisms-lemma-universally-closed-quasi-compact@cref}{{[originallemma][0][]04XU}{[1][0][]source}{}{}{}}
\makeatother
\crefname{originallemma}{original lemma}{original lemmas}
\Crefname{originallemma}{Original lemma}{Original lemmas}
\makeatletter
\newlabel{spaces-lemma-base-change-immersions}{{02YW}{source}{}{sourcepointer-4}{}}
\newlabel{spaces-lemma-base-change-immersions@cref}{{[originallemma][0][]02YW}{[1][0][]source}{}{}{}}
\makeatother
\crefname{originallemma}{original lemma}{original lemmas}
\Crefname{originallemma}{Original lemma}{Original lemmas}
\makeatletter
\newlabel{spaces-lemma-composition-immersions}{{02YV}{source}{}{sourcepointer-5}{}}
\newlabel{spaces-lemma-composition-immersions@cref}{{[originallemma][0][]02YV}{[1][0][]source}{}{}{}}
\makeatother
\crefname{originallemma}{original lemma}{original lemmas}
\Crefname{originallemma}{Original lemma}{Original lemmas}
\makeatletter
\newlabel{spaces-properties-lemma-points-cartesian}{{03H4}{source}{}{sourcepointer-6}{}}
\newlabel{spaces-properties-lemma-points-cartesian@cref}{{[originallemma][0][]03H4}{[1][0][]source}{}{}{}}
\makeatother
\crefname{originallemma}{original lemma}{original lemmas}
\Crefname{originallemma}{Original lemma}{Original lemmas}
\makeatletter
\newlabel{spaces-properties-lemma-reduced-closed-subspace}{{03IQ}{source}{}{sourcepointer-7}{}}
\newlabel{spaces-properties-lemma-reduced-closed-subspace@cref}{{[originallemma][0][]03IQ}{[1][0][]source}{}{}{}}
\makeatother
\title{A universally closed morphism of algebraic spaces over a scheme is quasi-compact}
\author{The Stacks Project Authors}
\date{stacks-project:spaces-morphisms.tex:local-snapshot-20261012; GNU FDL 1.2 or later}
\begin{document}
\maketitle
\setcounter{subsection}{1}
\noindent\textit{Editorial note.} All delivered mathematical text is the author's own, copied from the pinned Stacks Project chapter file spaces-morphisms.tex; no statement or proof step was rewritten, reordered or removed.\par\medskip
\noindent\textit{Editorial note.} The proof opens with the author's remark that it repeats the proof of the scheme case in Morphisms, Lemma morphisms-lemma-universally-closed-quasi-compact; that sentence is kept verbatim and the argument following it is written out in full, so nothing counted rests on the pointer.\par\medskip
\noindent\textit{Editorial note.} The seven references of the counted spans are recorded as non-load-bearing navigation pointers; the argument that consumes them stays inside the counted proof span.\par\medskip
\noindent\textit{Editorial note.} Preamble: the three native excerpt spans of sources/stacks-spaces-morphisms/preamble.tex (lines 8-13, the reference/slogan/history comment environments; 14-18, the Xy-pic setup; 154-217, the theorem declarations and mathematical macros). The chapter file itself carries no preamble text.\par\medskip
\noindent\textit{Allowed standard prerequisites (editorial).}\begin{itemize}
\item Quasi-compactness of a morphism of algebraic spaces and its criterion local on the target. Lemma lemma-quasi-compact-local, used at proof lines 1696-1700 and 1734-1737. Reduces the problem to the case Y = Spec(B) and converts the failure of quasi-compactness into a contradiction.
\item Permanence of quasi-compactness for the morphisms X\_i -> Y. Lemma lemma-quasi-compact-permanence, used at proof lines 1707-1708. Supplies the quasi-compactness of the pieces of the chosen open covering.
\item Reduced closed subspaces of an algebraic space cut out by a closed set of points. Properties of Spaces, Lemma spaces-properties-lemma-reduced-closed-subspace, used at proof lines 1710-1714. Produces the subspace Z whose image is shown not to be closed.
\item Points of fibre products of algebraic spaces and their field-valued representatives. Properties of Spaces, Lemma spaces-properties-lemma-points-cartesian, used at proof lines 1763-1765. Lifts the pair (x, t') to a point z of X x\_Y T lying in |Z|.
\item Base change of open immersions and of the diagram T x\_Y X -> Y. Spaces, Lemmas spaces-lemma-base-change-immersions and spaces-lemma-composition-immersions, used at proof lines 1715-1718. Identifies T\_i x\_Y X\_i as an open subspace so that Z is cut out by its complement.
\end{itemize}\medskip
\ifdefined\hypertarget\hypertarget{sourcepointer-1}{}\fi
\noindent\textit{Source reference.} Original lemma 03KG: focused-seeds/AGP/sources/stacks-snapshot-20261011/tags.txt, source lines 4676--4676. Navigation pointer into the Stacks Project; the standard fact it supplies is applied inside the counted proof span and the argument resting on it stays there.\par\medskip
\ifdefined\hypertarget\hypertarget{sourcepointer-2}{}\fi
\noindent\textit{Source reference.} Original lemma 03KS: focused-seeds/AGP/sources/stacks-snapshot-20261011/tags.txt, source lines 4688--4688. Navigation pointer into the Stacks Project; the standard fact it supplies is applied inside the counted proof span and the argument resting on it stays there.\par\medskip
\ifdefined\hypertarget\hypertarget{sourcepointer-3}{}\fi
\noindent\textit{Source reference.} Original lemma 04XU: focused-seeds/AGP/sources/stacks-snapshot-20261011/tags.txt, source lines 6407--6407. Navigation pointer into the Stacks Project; the standard fact it supplies is applied inside the counted proof span and the argument resting on it stays there.\par\medskip
\ifdefined\hypertarget\hypertarget{sourcepointer-4}{}\fi
\noindent\textit{Source reference.} Original lemma 02YW: focused-seeds/AGP/sources/stacks-snapshot-20261011/tags.txt, source lines 3891--3891. Navigation pointer into the Stacks Project; the standard fact it supplies is applied inside the counted proof span and the argument resting on it stays there.\par\medskip
\ifdefined\hypertarget\hypertarget{sourcepointer-5}{}\fi
\noindent\textit{Source reference.} Original lemma 02YV: focused-seeds/AGP/sources/stacks-snapshot-20261011/tags.txt, source lines 3890--3890. Navigation pointer into the Stacks Project; the standard fact it supplies is applied inside the counted proof span and the argument resting on it stays there.\par\medskip
\ifdefined\hypertarget\hypertarget{sourcepointer-6}{}\fi
\noindent\textit{Source reference.} Original lemma 03H4: focused-seeds/AGP/sources/stacks-snapshot-20261011/tags.txt, source lines 4553--4553. Navigation pointer into the Stacks Project; the standard fact it supplies is applied inside the counted proof span and the argument resting on it stays there.\par\medskip
\ifdefined\hypertarget\hypertarget{sourcepointer-7}{}\fi
\noindent\textit{Source reference.} Original lemma 03IQ: focused-seeds/AGP/sources/stacks-snapshot-20261011/tags.txt, source lines 4612--4612. Navigation pointer into the Stacks Project; the standard fact it supplies is applied inside the counted proof span and the argument resting on it stays there.\par\medskip
\begin{lemma}
\label{lemma-universally-closed-quasi-compact}
Let $S$ be a scheme.
A universally closed morphism of algebraic spaces over $S$ is quasi-compact.
\end{lemma}

\begin{proof}
This proof is a repeat of the proof in the case of schemes, see
Morphisms, Lemma \ref{morphisms-lemma-universally-closed-quasi-compact}.
Let $f : X \to Y$ be a morphism of algebraic spaces over $S$.
Assume that $f$ is not quasi-compact.
Our goal is to show that $f$ is not universally closed. By
Lemma \ref{lemma-quasi-compact-local}
there exists an affine scheme $Z$ and a morphism $Z \to Y$
such that $Z \times_Y X \to Z$ is not quasi-compact. To achieve our goal
it suffices to show that $Z \times_Y X \to Z$ is not universally closed,
hence we may assume that $Y = \Spec(B)$ for some ring $B$.

\medskip\noindent
Write $X = \bigcup_{i \in I} X_i$ where the $X_i$ are quasi-compact
open subspaces of $X$. For example, choose a surjective \'etale morphism
$U \to X$ where $U$ is a scheme, choose an affine open covering
$U = \bigcup U_i$ and let $X_i \subset X$ be the image of $U_i$.
We will use later that the morphisms $X_i \to Y$ are quasi-compact, see
Lemma \ref{lemma-quasi-compact-permanence}.
Let $T = \Spec(B[a_i ; i \in I])$. Let $T_i = D(a_i) \subset T$.
Let $Z \subset T \times_Y X$ be the reduced closed subspace whose underlying
closed set of points is
$|T \times_Y Z| \setminus \bigcup_{i \in I} |T_i \times_Y X_i|$, see
Properties of Spaces,
Lemma \ref{spaces-properties-lemma-reduced-closed-subspace}.
(Note that $T_i \times_Y X_i$ is an open subspace of $T \times_Y X$ as
$T_i \to T$ and $X_i \to X$ are open immersions, see
Spaces, Lemmas \ref{spaces-lemma-base-change-immersions} and
\ref{spaces-lemma-composition-immersions}.) Here is a diagram
$$
\xymatrix{
Z \ar[r] \ar[rd] &
T \times_Y X \ar[d]^{f_T} \ar[r]_q &
X \ar[d]^f \\
& T \ar[r]^p & Y
}
$$
It suffices to prove that the image $f_T(|Z|)$ is not closed in $|T|$.

\medskip\noindent
We claim there exists a point $y \in Y$ such that there is no
affine open neighborhood $V$ of $y$ in $Y$ such that $X_V$ is quasi-compact.
If not then we can cover $Y$ with finitely many such $V$ and for
each $V$ the morphism $Y_V \to V$ is quasi-compact by
Lemma \ref{lemma-quasi-compact-permanence}
and then
Lemma \ref{lemma-quasi-compact-local}
implies $f$ quasi-compact, a contradiction. Fix a $y \in Y$ as in the claim.

\medskip\noindent
Let $t \in T$ be the point lying over $y$ with $\kappa(t) = \kappa(y)$
such that $a_i = 1$ in $\kappa(t)$ for all $i$. Suppose $z \in |Z|$ with
$f_T(z) = t$. Then $q(t) \in X_i$ for some $i$. Hence $f_T(z) \not \in T_i$
by construction of $Z$, which contradicts the fact that $t \in T_i$ by
construction. Hence we see that $t \in |T| \setminus f_T(|Z|)$.

\medskip\noindent
Assume $f_T(|Z|)$ is closed in $|T|$. Then there exists an element
$g \in B[a_i; i \in I]$ with $f_T(|Z|) \subset V(g)$ but $t \not \in V(g)$.
Hence the image of $g$ in $\kappa(t)$ is nonzero. In particular some
coefficient of $g$ has nonzero image in $\kappa(y)$. Hence this coefficient is
invertible on some affine open neighborhood $V$ of $y$. Let $J$ be the finite
set of $j \in I$ such that the variable $a_j$ appears in $g$.
Since $X_V$ is not quasi-compact and each $X_{i, V}$ is quasi-compact,
we may choose a point $x \in |X_V| \setminus \bigcup_{j \in J} |X_{j, V}|$.
In other words, $x \in |X| \setminus \bigcup_{j \in J} |X_j|$ and $x$ lies
above some $v \in V$. Since $g$ has a coefficient that is invertible on $V$,
we can find a point $t' \in T$ lying above $v$ such that $t' \not \in V(g)$ and
$t' \in V(a_i)$ for all $i \notin J$. This is true because
$V(a_i; i \in I \setminus J) = \Spec(B[a_j; j\in J])$
and the set of points of this scheme lying over $v$ is bijective
with $\Spec(\kappa(v)[a_j; j \in J])$ and $g$ restricts to
a nonzero element of this polynomial ring by construction.
In other words $t' \not \in T_i$ for each $i \not \in J$. By
Properties of Spaces, Lemma \ref{spaces-properties-lemma-points-cartesian}
we can find a point $z$ of $X \times_Y T$ mapping to $x \in X$ and to
$t' \in T$. Since $x \not \in |X_j|$ for $j \in J$ and $t' \not \in T_i$
for $i \in I \setminus J$ we see that $z \in |Z|$. On the other hand
$f_T(z) = t' \not \in V(g)$ which contradicts $f_T(Z) \subset V(g)$.
Thus the assumption ``$f_T(|Z|)$ closed'' is wrong and we conclude indeed
that $f_T$ is not closed as desired.
\end{proof}

\end{document}
